\documentclass{article}
\usepackage[UTF8]{ctex} % 使用 ctex 宏包
\usepackage{graphicx}%插入和管理图片（如 .png, .jpg, .pdf 等）
\usepackage{caption}     % 图片标题
\usepackage{array}%增强表格功能，自定义列格式和样式
\usepackage{fancyhdr}%自定义页眉（Header）和页脚（Footer）
\usepackage{lastpage}%获取文档总页数
\usepackage{geometry}%调整页面布局参数（如边距、纸张大小)
\usepackage{fontspec}  % XeLaTeX专用字体包
\usepackage{amsmath} % 数学公式支持
\usepackage{amsfonts} % 数学字体支持
\usepackage{ulem}%下划线包
\usepackage{gensymb}%添加度的符号
\usepackage{multirow}%合并表格的行
\usepackage{siunitx}%添加度分秒的符号
\usepackage{etoolbox}%调整表格等与文字的间距


% 调整浮动体与周围内容的间距
\makeatletter
\setlength{\textfloatsep}{5pt plus 2pt minus 4pt} % 与上下正文间距
\setlength{\floatsep}{5pt plus 2pt minus 4pt}      % 浮动体之间的间距
\setlength{\intextsep}{5pt plus 2pt minus 4pt}     % 页面内浮动体与文本的间距
\makeatother

\setlength{\parindent}{0pt} % 取消所有段落缩进


% 设置中文字体（默认宋体）
\setmainfont{SimSun}   % 宋体
\setsansfont{Microsoft YaHei}  % 如果需要微软雅黑等其他字体

\geometry{left=2cm,right=2cm,top=2.5cm,bottom=2.5cm}

% 设置页眉页脚
\pagestyle{fancy}
\fancyhf{}
\renewcommand{\headrulewidth}{0pt}
\fancyfoot[C]{浙江大学物理实验教学中心制 \quad 第 \thepage 页 / 共 \pageref{LastPage} 页}

%正文开始
\begin{document}

\linespread{1.5}  % 设置为1.5倍行距
\selectfont

\begin{center}
\vspace*{1em}
\includegraphics[width=6.2cm]{assets/zju-name.jpg}\\
\vspace{0.6cm}
{\fontsize{50pt}{18pt}\selectfont \textbf{物理学实验I}}\\
\vspace{2em}
{\fontsize{50pt}{18pt}\selectfont \textbf{电子实验报告}}\\
\vspace{10em}

\Large \textbf{实验名称：\uline{\makebox[6.5cm][c]{}}}\\
\Large \textbf{实验桌号：\uline{\makebox[6.5cm][c]{}}}\\
\Large \textbf{指导教师：\uline{\makebox[6.5cm][c]{}}}\\
\vspace{4em}
\Large \textbf{班 \quad 级：\uline{\makebox[6cm][c]{}}}\\
\Large \textbf{姓 \quad 名：\uline{\makebox[6cm][c]{}}}\\
\Large \textbf{学 \quad 号：\uline{\makebox[6cm][c]{}}}\\
\vspace{4em}
\Large \textbf{实验日期： 年月日 星期下午}\\
\end{center}

\newpage

% 主要内容
\section*{一、 预习报告}
\textit{（注：将已经写好的“电子预习报告”内容拷贝过来）} 


\subsection*{1. 实验综述}
\textit{（自述实验现象、实验原理和实验方法，不超过300字，5分）} \\


\subsection*{2. 实验重点}
\textit{（简述本实验的学习重点，不超过100字，3分）} \\


\subsection*{3. 实验难点}
\textit{（简述本实验的实现难点，不超过100字，2分）} \\


\section*{二、 原始数据}
\textit{（将有老师签名的“自备数据记录草稿纸”的扫描或手机拍摄图粘贴在下方,20分）} \\


\section*{三、 结果与分析}
\subsection*{1. 数据处理与结果}
\textit{（列出数据表格、选择数据处理方法、给定测量或计算结果，20分）} \\


\subsection*{2. 误差分析}
\textit{（运用测量误差、相对误差、不确定度等分析实验结果，20分）} \\


\subsection*{3. 实验探讨}
\textit{（对实验内容、现象和过程的小结，不超过100字，10分）} \\


\section*{四、 思考题}
\textit{（解答课本或讲义或老师布置的思考题，20分）} \\


\end{document}
